% ---------------------------------------------------------------------------
\section{Proofs}
\label{app:proofs}
% ---------------------------------------------------------------------------

This appendix proves the results of Sec.~\ref{sec:theory} and gives the details of the
ring-lattice windows, of the fusion of rings through shared pairs and of the cost
estimate.

\subsection{Determinism, locality and nesting}

\begin{proof}[Proof of Proposition~\ref{prop:determinism}]
$\mathcal T_t$ is defined pointwise from $G$ and $t$, the adjacency of $L_t$ pointwise
from $\mathcal T_t$, and the connected components of a graph are unique.
\end{proof}

\begin{proof}[Proof of Proposition~\ref{prop:local}]
Two triples are adjacent in $L_t$ iff they share a pair $\{x,y\}$. If $\{x,y,w\}$ is a
wedge its centre is $x$, $y$ or $w$, and in each case $w$ is adjacent to $x$ or $y$.
Every neighbour of a visited triple is therefore generated from one of its three pairs,
so the search traverses the component and nothing else. Admissibility of
$\{x,y,w\}$ depends on $N[x],N[y],N[w]$ only.
\end{proof}

\begin{proof}[Proof of Theorem~\ref{thm:nesting}]
$J\ge t_2$ implies $J\ge t_1$ and the wedge condition does not involve $t$, so
$\mathcal T_{t_2}\subseteq\mathcal T_{t_1}$. Adjacency (sharing two nodes) does not
involve $t$ either, so $L_{t_2}$ is induced. A connected subgraph of $L_{t_2}$ is
connected in $L_{t_1}$, so each component of $L_{t_2}$, and the union of its triples,
lies inside a single component of $L_{t_1}$.
\end{proof}

\subsection{The two ends of the hierarchy}

\begin{proof}[Proof of Lemma~\ref{lem:constant}]
$\pi_{\mathcal C}$ is concentrated on $c$, so $\Omega_e=\pi_{\mathcal G}(c)<1$; a pair agrees
exactly when $c_{\mathcal G}(i,j)=c$, so $\Omega_u=\pi_{\mathcal G}(c)$ as well.
\end{proof}

\begin{proof}[Proof of Proposition~\ref{prop:endpoints}]
(i) At $t=0$ every wedge is admissible. Two edges sharing an endpoint form a wedge, and
the wedges built on a common edge share its endpoints, so a walk in the line graph of
$G$ is a walk in $L_0$. The line graph of a connected graph is connected, and every
node of a connected graph with $|V|\ge3$ lies on a wedge. The single community $V$
gives $c_{\mathcal C}\equiv1$, and Lemma~\ref{lem:constant} applies.
(ii) $J(u,v)=1$ iff $N[u]=N[v]$, an equivalence relation whose classes are cliques; a
triple is admissible at $t=1$ iff it lies within one class, and triples within one
class are connected in $L_1$. Without classes of size at least three, $c_{\mathcal C}\equiv0$
and Lemma~\ref{lem:constant} applies. For the final claim: the cover changes only at
the finitely many distinct triple scores, so $\Omega(t)$ takes finitely many values on
$[0,1]$ and its maximum $\Omega^\star$ is attained; under the stated hypotheses
$\Omega(0)=\Omega(1)=0$ by (i) and (ii), so since $\Omega^\star\ge\Omega(t)>0$ every
maximiser lies in $(0,1)$.
\end{proof}

\subsection{Ring lattices}

\begin{proof}[Proof of Lemma~\ref{lem:ringJ}]
$N[x]$ is an arc of $k+1$ consecutive ring nodes, and two arcs at offset $d\le k$
overlap in $k+1-d$ nodes when $s\ge 2k+2$, which gives $J(d)$. A merged node's closed
neighbourhood is the union of its two arcs, $2k+1$ nodes, and the arc of $x$ meets
only the one in its own ring, which gives $J_u(d)$. Nodes on different rings can share
only merged nodes, and isolation leaves at most one, $u$, in both neighbourhoods.
\end{proof}

\begin{proof}[Proof of Theorem~\ref{thm:window}]
\emph{No merges.} $J$ decreases with ring distance, and any three distinct ring nodes
include a pair at distance at least two, so no triple is admissible when $t>J(2)$ and
every ring is lost. For $t\le J(2)$ the consecutive triples $\{i,i{+}1,i{+}2\}$, with
pair distances $1,1,2$, are admissible and each shares a pair with the next, so they
form a single component covering the ring. Every other admissible wedge lies in a
cyclic interval of length at most $k$, and moving its outer nodes inward preserves
admissibility and connects it to a consecutive triple, so there are no further
components. No wedge spans two rings, which share no edge, so the components are
exactly the rings.

\emph{Lower edge.} A wedge with nodes on two rings needs a node adjacent to both, i.e.
a merged node $u$; its weakest pair is $J(z,y')=1/(2k+1)$ by Eq.~\eqref{eq:ringJ}, the
other two being of the form $J_u(d)$ with $d\le k/2$, which is larger. If
$t\le1/(2k+1)$, the triple $\{z,u,y'\}$ is admissible and, since $J_u(1)$ and $J(2)$
exceed $1/(2k+1)$ for $k\ge4$, so are triples of both rings through the pairs
$\{z,u\}$ and $\{u,y'\}$: the two rings fall into one component. If $t>1/(2k+1)$ no
admissible triple contains nodes of two rings other than $u$, so each component lies
within one ground-truth community.

\emph{Upper edge.} Above the lower edge every admissible triple lies within one ring,
and each wedge lies in an arc of length at most $k$, which by isolation contains at
most one merge position. Fix a merged node $u$ at position $0$ of one of its rings.
Every admissible triple through $u$ on this ring has two $u$-pairs, at distinct
distances or on opposite sides of $u$, so the best is $T_u=\{-1,0,1\}$, with weakest
pairs $J_u(1)=k/(2k+2)$: for $t>J_u(1)$, $u$ is in no triple of the ring and is missing
from its community. For $J_u(2)<t\le J_u(1)$, $T_u$ is the only admissible triple through
$u$ on this ring; the corresponding triple on $u$'s other ring shares only $u$ with it
and cannot attach it, so $T_u$ connects to the ring only through its pair $\{-1,1\}$.
Every completion of that pair other than $u$ has score at most $J(3)=(k-2)/(k+4)$,
attained by $\pm2$. Hence for $\max(J(3),J_u(2))<t\le J_u(1)$, $T_u$ is an isolated
component and $\{-1,0,1\}$ appears as a spurious community.

Conversely, let $t\le J_u(1)$ and $t\le\max(J(3),J_u(2))$. The consecutive triples
that avoid $u$ are admissible (their pairs are at distances $1$, $1$, $2$, and
$J(2)>J_u(1)\ge t$ for $k\ge4$), and so is $T_u$. If $t\le J(3)$, the bridges
$\{-2,-1,1\}$ and $\{-1,1,2\}$, whose pairs avoid $u$ and lie at distances $1$, $2$ and
$3$, are admissible; they share $\{-2,-1\}$ and $\{1,2\}$ with the consecutive triples
on either side of the merge and $\{-1,1\}$ with $T_u$. If instead $t\le J_u(2)$, the
consecutive triples $\{-2,-1,0\}$ and $\{0,1,2\}$ through $u$ are admissible and share
$\{-1,0\}$ and $\{0,1\}$ with $T_u$. Either way the triples listed so far form one
component per ring. Every other admissible triple $T$ of the ring joins it through
\emph{contraction moves}: replace an outer node of $T$ whose next node towards the
other two (skipping $u$ if $u\notin T$) is not already in $T$ by that node, and stop
when the three nodes are consecutive in the ring order (with $u$ removed if
$u\notin T$). A move keeps two nodes of $T$, so successive triples share a pair; it
increases no ring distance within the triple, so the triple remains a wedge; and it
does not decrease the score, because pairs not involving $u$ have commonality $J(d)$
and pairs involving $u$ have $J_u(d)$, both decreasing in $d$, and a move that replaces
$u$ itself turns each $u$-pair at distance $d$ into a pair at distance $d'\le d$ with
$J(d')\ge J(d)>J_u(d)$. The moves therefore
end, through admissible triples, at a consecutive triple, a bridge, or a triple
through $u$ already listed, so all admissible triples within each ring are connected.
Since $J_u(1)<J(2)$ for $k\ge4$, the no-merge condition is implied, and
Eq.~\eqref{eq:window} follows.
\end{proof}

\emph{Isolation.} The proof uses isolation only through the facts that an arc of $k+1$
consecutive nodes contains at most one merge position and that nodes of different rings
share at most one merged neighbour. For a ring pair merged twice the separation $k+1$ is
sharp: shared nodes $k$ apart on both rings raise the lower edge to $1/k$.

\emph{Unequal degrees.} For one isolated merge between rings of even degrees
$k_a,k_b\ge4$ and sizes at least $2k_a+2$ and $2k_b+2$, the same argument applies with
the commonalities of Lemma~\ref{lem:ringJ} for unequal degrees. The cross pair has
$J=1/(k_a+k_b+1)$, which gives the lower edge of Eq.~\eqref{eq:windowab}. On the ring of
degree $x$, with $y$ the degree of the other ring, the upper-edge argument runs with
$J_u(1)=x/(x+y+2)$, $J_u(2)=(x-1)/(x+y+3)$ and $J(3)=(x-2)/(x+4)$, which gives
$U(x,y)$; both rings must be recovered, which gives the minimum over $U(k_a,k_b)$ and
$U(k_b,k_a)$.

\emph{Obstruction for unequal degrees.} A shared node between a dense and a sparse
community is held more weakly by the sparse one, so the window of such a merge narrows
and moves down, with upper edge $k_b/(k_a+k_b+2)$ for $k_a\gg k_b$. It cannot overlap
the window of merges between two sparse rings, whose lower edge is $1/(2k_b+1)$, once
$k_a\ge2k_b^2-2$: beyond that degree contrast no single threshold recovers a family that
also contains merges between two sparse rings, even without noise.

\subsection{Shared pairs and pair-disjointness}

\begin{proof}[Proof of Proposition~\ref{prop:sharedpairs}]
If $T_A\ne T_B$, the two triples share exactly the two nodes $u_1,u_2$ and are adjacent
in $L_t$.
\end{proof}

\emph{Fusion of rings.} Connectivity through shared pairs is what lets a community grow
through its own open wedges, and Proposition~\ref{prop:sharedpairs} is its counterpart.
On ring lattices, three communities that pairwise share single
nodes are joined when two of the three shared-node pairs are adjacent and the third lies
within distance $k$; for $k=8$ and pairwise adjacent shared nodes two apart, the rings
are recovered only for $4/13<t\le4/9$, and for $k=4$ they are never recovered. Two rings
joined only through two shared merged nodes, with no other merge within ring distance
$k$ of them, can be fused above the lower edge $1/(2k+1)$ only when those nodes lie
within ring distance $k$ of each other on both rings, or within $k/2$ on one ring and
within $3k/2$ on the other; for $k=8$, shared nodes $2$ apart on one ring and $11$ apart
on the other fuse the rings for $t\le3/23$. Theorem~\ref{thm:window}'s isolation
hypothesis excludes both cases. Whether fusion then occurs depends on the threshold,
since the merge lowers the commonality of the shared pair: two rings with $s=50$, $k=8$
that share two nodes four positions apart are fused at small $t$ but recovered
separately at $t=0.4$, where the shared pair, with $J=1/3$, is not admitted.

\begin{proof}[Proof of Corollary~\ref{cor:pairdisjoint}]
Two distinct admissible triples that contain the same pair share exactly two nodes and
are therefore adjacent in $L_t$, so all admissible triples through the pair lie in one
component.
\end{proof}

\subsection{Cost}

Every commonality needed is available from the wedges themselves: for a pair $\{a,b\}$,
$|N[a]\cap N[b]|$ is the number of common neighbours of $a$ and $b$, that is of centred
wedges $a$--$c$--$b$ (a triangle contributes once to each of its three pairs), plus $2$
if $a\sim b$, so one pass over the centred wedges, with pairs indexed by hashing, counts
all of them. Our implementation finds components on the incidence graph between admitted
triples and their pairs over a grid of thresholds. Alternatively, as in sequential clique
percolation~\cite{kumpula2008sequential}, a single union--find sweep over the triples in
decreasing order of score builds the merge tree of the triple components at all
thresholds in $O(W\log W)$, from which the cover at any one threshold is read in $O(W)$
more. For heavy-tailed degree distributions the hubs dominate through $W$, as for
every triangle-based method.

\emph{Degree-ratio pruning.} In the pruning of Sec.~\ref{sec:cost}, every wedge through
an edge with $\rho<t_{\min}$ is discarded. Common neighbours reached through discarded
edges must still be counted, for example by intersecting the adjacency lists of the kept
pairs; counting only the kept wedges would underestimate $J$. This is the size filter of
exact set-similarity joins~\cite{arasu2006efficient,bayardo2007scaling}, applied to
closed neighbourhoods of sizes $d_u+1$ and $d_v+1$. The pruning applies to every rule
that tests the edges of a wedge, including our three ablations, but not to link
communities, which score a wedge by its non-adjacent pair alone, whose commonality is
not bounded by the degrees of the wedge's edges, so every wedge through a hub has to be
enumerated. Our implementation counts the common neighbours of the kept pairs from the
full adjacency matrix, and pruned and unpruned runs give identical covers on its test
graphs, which include a hub joined to many leaves, and on the full Amazon and DBLP graphs
at all 30 thresholds $t\ge0.02$ of our grid.
